% ECONOMICS_PROBLEM: {"id": "D83-586", "title": "Trust in recommendations", "jel_code": "D83", "source_id": "OP-0404"}
% !TeX program = xelatex
\documentclass[11pt,a4paper]{article}
\usepackage[margin=23mm]{geometry}
\usepackage{fontspec}
\setmainfont{Latin Modern Roman}
\usepackage{amsmath,amssymb,mathtools}
\usepackage{xcolor,enumitem,booktabs,longtable,array,xurl,hyperref}
\definecolor{muted}{HTML}{53636C}
\hypersetup{colorlinks=true,urlcolor=blue,linkcolor=blue}
\setlength{\parindent}{0pt}
\setlength{\parskip}{5pt}
\newcommand{\R}{\mathbb R}
\newcommand{\E}{\mathbb E}
\newcommand{\N}{\mathbb N}
\newcommand{\ind}{\mathbf 1}
\newcommand{\Var}{\operatorname{Var}}
\newcommand{\Cov}{\operatorname{Cov}}
\newcommand{\supp}{\operatorname{supp}}
\DeclareMathOperator*{\argmin}{arg\,min}
\DeclareMathOperator*{\argmax}{arg\,max}
\newcommand{\entrymeta}[3]{{\sffamily\small\color{muted}\textbf{\texttt{#1}}\quad|\quad #2\par #3\par}}
\newcommand{\Problemref}[1]{\texttt{#1}}
\newcommand{\JELref}[2]{\texttt{#2} (JEL \texttt{#1})}
\begin{document}
\section*{Trust in recommendations}
\textbf{Problem ID:} \texttt{D83-586}\quad \textbf{JEL:} \texttt{D83}\par
% BEGIN ECONOMICS STATEMENT
\label{problem:OP-0404}
{\sffamily\small\color{muted}Original subject group: Digital, platform and data economics\par}
\label{definitions:problem:30007686}
\textbf{Audit classification:} Published research agenda. The recommendation-trust dilemma and transparency agenda are verified. Additional manipulation loss was previously undefined and not necessarily nonnegative.
\subsection*{Definitions and precise research target}
Consider users receiving purchase recommendations from a retail platform over repeated visits. The platform can disclose a transparency message \(z\in\mathcal Z\), where \(\mathcal Z\) is a prespecified finite set of explanations about recommendation criteria, sponsorship and uncertainty. Let \(\theta_i\) denote user \(i\)'s stable preferences, \(b_i(z)\) the user's belief about recommendation quality after message \(z\), and \(a_i(z)\) the chosen product or outside option. Let \(v_i(a)\) be independently assessed money-valued utility net of price for action \(a\), and let \(a_i^*\) maximize that utility among feasible actions. Under audited, nonmanipulated seller information, define decision loss \(L(z)=E[v_i(a_i^*)-v_i(a_i(z))]\), expectation being over registered users in the study. Also define \(M(z)\) as the change in expected decision loss between the strategic-seller and audited environments under that disclosure. The design target is to minimize \(L(z)+M(z)\), subject to truthful disclosure and a specified implementation-cost limit. Randomize explanations and recommendation quality independently where feasible, eliciting beliefs separately from preferences. This permits distinguishing warranted rejection from distrust of useful advice. Study persistence after repeated experience and seller adaptation. The problem concerns appropriate reliance, so increasing recommendation acceptance is not itself the objective and complete algorithm disclosure is only one candidate design.

\textbf{Definitions, assumptions and mathematical qualifications.} Let \(a_i^*\) be a measurable utility maximizer among a specified finite action set, with ties resolved in advance. Define audited and strategic environments \(A,S\), with losses \(L_A(z)=E[v_i(a_i^*)-v_i(a_i^A(z))]\) and \(L_S(z)=E[v_i(a_i^*)-v_i(a_i^S(z))]\). Set \(M(z)=L_S(z)-L_A(z)\), which may be negative unless an additional restriction excludes beneficial strategic responses. Then \(L(z)+M(z)=L_S(z)\); the decomposition must use the same utility and population. Messages must be truthful under a verification standard and meet a stated cost bound. Belief elicitation does not identify stable preferences without a separate measurement model.
\subsection*{Variants and assumption changes}
\begin{enumerate}
\item \textbf{Editorial, no strategic response.} In a factorial experiment vary explanation and audited recommendation quality independently; identify appropriate-reliance losses under fixed sellers.
\item \textbf{Editorial, repeated adaptation.} Add strategic sellers and learning users, fix a finite response horizon and equilibrium selection, and minimize total strategic-environment decision loss over feasible truthful messages.
\end{enumerate}
\subsection*{References and source audit}
\textbf{Checked source.} 2025 conference draft, Section 1.1, fifth research direction, PDF p. 5. Full primary text retrieved. \url{https://conference.nber.org/confer/2025/DTs25/farronato.pdf}.
\begin{enumerate}\raggedright
\item\hypertarget{definitions:source:30007686:1}{} Chiara Farronato. 2025. \emph{Digital Platforms as Information Aggregators: Implications for their Design and Regulation}. 2025 conference draft, Section 1.1, fifth research direction, PDF p. 5.
\url{https://conference.nber.org/confer/2025/DTs25/farronato.pdf}.
\end{enumerate}

\subsection*{Common conventions: Source mathematical conventions}

These conventions apply to each entry unless explicitly replaced.

\textbf{Models and identification.} A primitive model \(m\) specifies all probability laws, feasible choices, information, equilibrium and measurement rules used by its target. The admissible class \(\mathcal M\) is fixed independently of outcomes. If \(P_O(m)\) is its observable probability law and \(\tau(m)\) the target, its identified set at observable law \(P\) is
\[
 \mathcal I_\tau(P)=\{\tau(m):m\in\mathcal M,\ P_O(m)=P\}.
\]
Point identification means that this set is a singleton. An empty set signals incompatibility of data and assumptions. Coordinate bounds need not describe the joint set. Bounds are sharp only if every retained target is attainable or arbitrarily approachable by compatible models. Finite bounds require explicit restrictions on unobserved quantities; unrestricted confounding, hidden wealth, extrapolation or arbitrary equilibrium selection can yield uninformative sets.

\textbf{Causal interventions.} A potential outcome \(Y_i(p)\) is the outcome under a complete policy regime \(p\), including timing, financing, information and market spillovers. The target population and units are fixed before treatment. With interference, use the complete assignment vector or a justified exposure mapping; individual no-interference notation alone cannot identify an economy-wide intervention. A difference of mean potential outcomes does not require the cross-world joint distribution. Conditioning on post-treatment migration, survival, enrollment or employment changes the estimand and needs separate justification.

\textbf{Statistical statements.} An estimator, confidence set, forecast or test specifies a sampling law, model class and loss/coverage criterion. Uniform \((1-\alpha)\)-coverage means \(\inf_{m\in\mathcal M}\Pr_m\{\tau(m)\in C_n\}\geq1-\alpha\), or an explicitly asymptotic version. The whole parameter space trivially has coverage, so informativeness requires a stated width, risk or power objective. A forecasting improvement is relative to a named benchmark, information set and evaluation population.

\textbf{Equilibrium and optimization.} An equilibrium comprises feasible plans, optimality, beliefs where needed, budget constraints and all market/resource clearing. The primitives or a stated benchmark define each correspondence \(\mathcal E(p;m)\). Nonemptiness is assumed or proved, never inferred just from compact policy sets. A selected-equilibrium target uses a declared selection rule; a robust target quantifies over all permitted equilibria. Use suprema or infima unless attainment is established. An \(\varepsilon\)-optimal policy is feasible and within \(\varepsilon\) of the finite optimum value. Report an empty feasible set separately.

\textbf{Welfare and accounting.} Normative utility, interpersonal normalization, social weights, numeraire, discounting and the use of public receipts are primitives. Behavioral utility need not coincide with the welfare criterion. Household welfare includes ownership income and rents once; adding profits again to compensating variation generally double-counts them. A monetary equivalent is defined by an explicit transfer and price convention, with monotonicity and range conditions. Average effects, mechanism contributions, transfers and welfare losses are different targets. Infinite sums require convergence and dynamic feasible plans require terminal or no-Ponzi conditions.

\textbf{Source versions.} A working-paper date, revision date and journal publication year are distinguished. A DOI for a journal version is not automatically the DOI of an earlier manuscript. Locators refer to the stated version and printed pages where specified. A metadata-only check does not verify a claim about a full-text conclusion. Every entry's source subsection narrows the provenance claim accordingly.
% END ECONOMICS STATEMENT
\end{document}
